\documentclass[10pt,a4paper]{amsart}
\usepackage[margin=19mm]{geometry}
\usepackage{amsmath,amssymb,mathtools}
\usepackage[scr=rsfs,cal=euler]{mathalfa}
\usepackage{xfrac}
\usepackage[unicode,hidelinks,bookmarks=false]{hyperref}
\newcommand{\ie}{i.e.\ }
\newcommand{\cf}{cf.\ }
\newcommand{\resp}{respectively\ }
\newcommand{\Subsection}{\subsection}
\providecommand{\nobreakdash}{\nobreak\hbox{-}}
\setlength{\emergencystretch}{2em}
\multlinegap0pt
\usepackage{bm}
\usepackage{bbm}
\usepackage{dsfont}
\usepackage{xspace}
\usepackage{cancel}
\usepackage{mathtools}
\usepackage{amsthm}
\makeatletter
\@ifundefined{thm}{\newtheorem{thm}{Thm}}{}
\@ifundefined{lem}{\newtheorem{lem}[thm]{Lemma}}{}
\makeatother
\title[Chern character and Fermi point]{Chern character and Fermi point}
\author{Kyouhei Horie}
\date{Source: arXiv:2505.06704v2 (CC BY 4.0)}
\begin{document}
\maketitle

\noindent\emph{Source and licence.} Chern character and Fermi point. Version arXiv:2505.06704v2 (CC BY 4.0), distributed under the Creative Commons Attribution 4.0 International licence, as recorded in the arXiv licence metadata for this exact version and shown on the abstract page https://arxiv.org/abs/2505.06704v2. Full rights notice: licenses/arxiv-2505.06704-CC-BY-4.0.txt.\par\smallskip

\noindent\textit{Editorial scope.} Counted unit: the Lem whose source label is \texttt{lem:6.5} in the article (source0.tex lines 3008--3023, source label \texttt{lem:6.5}), delivered together with the article's own complete original proof of it (source0.tex lines 3026--3122, one \texttt{proof} environment of 97 lines). No mathematics is written, repaired or re-derived: statement and proof are the article's own text line for line. Required context retained uncounted: none. Every definition, display and label of those lines is retained unchanged. Omitted from this package and not redistributed: 1--3007, 3024--3025, 3123--3480. Nothing inside a retained excerpt is omitted or altered: every retained body is delivered exactly as the article's lines are written, so no omission receipt of any kind accompanies this package. Citations: the retained excerpts carry no citation command at all, so no bibliography entry of the article is transcribed and no reference is redistributed. Reference adaptation: none was needed and none was made; every reference inside the delivered unit resolves inside it, so no unresolved reference or citation is produced. Printed slips kept literal: no printed word, symbol, hypothesis or number of the counted statement or of its proof was altered. Layout and class: the article's own class is \texttt{amsart} with packages including amsmath, amssymb, amsthm, graphicx, amscd, tikz, tikz-cd, hyperref; the wrapper is the lane's pinned \texttt{amsart} editorial wrapper at 10pt on A4, which restates the theorem-like environments the retained text needs and adds the article's own macro definitions that the retained text needs (none), with extra wrapper packages (bm, bbm, dsfont, xspace, cancel, mathtools). The delivered numbering is the unit's own. Assets: no retained span contains a figure environment, a graphics-inclusion command, a PDF-page inclusion, a table or a TikZ picture; no image file is copied into this package, the package declares no asset dependency and no image file is redistributed with it. Licence: version arXiv:2505.06704v2 of this article is distributed under the Creative Commons Attribution 4.0 International licence, as recorded in the arXiv licence metadata for this exact version and shown on the abstract page https://arxiv.org/abs/2505.06704v2. A full rights notice is delivered at licenses/arxiv-2505.06704-CC-BY-4.0.txt. Characters: every retained excerpt is pure ASCII, so the delivered engine, XeLaTeX with its default Latin Modern text font, reports no missing character.
\begin{lem}\label{lem:6.5}
We suppose that the conditions \(c \neq 0\), \(b = 0\), and \(a = E\) hold. Under these assumptions, we have \(\Delta > 0\) if and only if \(|c| \neq 1\). Furthermore, the equation \((H^\# - E)\psi = 0\) has nontrivial \(l^2\)-solutions if and only if \(|c| < 1\). Additionally, when the condition \(|c| < 1\) is satisfied, we have the following properties of the kernel:
\begin{align*}
\ker(H^\# - E) &=
\begin{cases}
\langle e_1 \rangle & a \neq 0, \\
\langle e_1, e_4 \rangle & a = 0,
\end{cases}
&
\dim(\ker(H^\# - E)) &=
\begin{cases}
1 & a \neq 0, \\
2 & a = 0.
\end{cases}
\end{align*}
\end{lem}

\begin{proof}
We have \(\Delta = (|c|^2 - 1)^2\), which is greater than zero if and only if the norm of \(c\) is not equal to 1. Thus, we proceed under the assumption that \(|c| \neq 1\) holds throughout the following discussion.

\begin{enumerate}
    \item The condition \(|c| < 1\) is satisfied. In this case, we find that the eigenvalues are \(\lambda_1 = c\) and \(\lambda_2 = \frac{1}{\bar{c}}\). We introduce a matrix \(T\) defined as follows:

    \begin{align*}
    T = \begin{pmatrix}
    0 & 1 & 0 & \frac{2a \bar{c}}{|c|^2 - 1} \\
    0 & 0 & 0 & 1 \\
    -\frac{2a c}{|c|^2 - 1} & 0 & 1 & 0 \\
    1 & 0 & 0 & 0
    \end{pmatrix}.
    \end{align*}

    Since the determinant of \(T\) is $1$, we conclude that \(T\) is an invertible matrix. The inverse is

    \begin{align*}
    T^{-1} = \begin{pmatrix}
    0 & 0 & 0 & 1 \\
    1 & -\frac{2a \bar{c}}{|c|^2 - 1} & 0 & 0 \\
    0 & 0 & 1 & \frac{2a c}{|c|^2 - 1} \\
    0 & 1 & 0 & 0
    \end{pmatrix}.
    \end{align*}

    Under this transformation, we verify that \(T^{-1} R T = \text{diag}(c, c, \frac{1}{\bar{c}}, \frac{1}{\bar{c}})\). We express the transformed solution \(T^{-1} \psi(n)\) as the transpose \((x(n), y(n), z(n), w(n))\). For \(n = 1\), we obtain:

    \begin{align*}
    T^{-1} \psi(1) = z \begin{pmatrix} 0 \\ \bar{c} \\ 0 \\ 0 \end{pmatrix} + w \begin{pmatrix} \bar{c} \\ 0 \\ \frac{2a}{|c|^2 - 1} \\ 0 \end{pmatrix},
    \end{align*}

    and we compute the \(l^2\)-norm of \(T^{-1} \psi\) as:

    \begin{align*}
    \| T^{-1} \psi \|^2 = \sum_{n \in \mathbb{N}} \left( |c|^{2n} \big( |x(n)|^2 + |y(n)|^2 \big) + \left| \frac{1}{\bar{c}} \right|^{2n} \big( |z(n)|^2 + |w(n)|^2 \big) \right).
    \end{align*}

    Since \(|\lambda_1| |\lambda_2| = 1\) and \(|c| < 1\), we conclude that \(\| T^{-1} \psi \|^2 < \infty\) holds if and only if either \(z(1) = w(1) = 0\) or both \(x(1) = y(1) = 0\) and \(\left| \frac{1}{\bar{c}} \right| < 1\).

    \begin{enumerate}
        \item We suppose that \(z(1) = w(1) = 0\) holds. Since \(w(1) = w \frac{2a}{|c|^2 - 1} = 0\), we deduce that either \(w = 0\) or \(a = 0\). Therefore, we determine that the kernel of \(H^\# - E\) is:

        \begin{align*}
        \ker(H^\# - E) =
        \begin{cases}
        \langle e_1 \rangle & \text{if } a \neq 0, \\
        \langle e_1, e_4 \rangle & \text{if } a = 0.
        \end{cases}
        \end{align*}

        \item We suppose that \(\left| \frac{1}{\bar{c}} \right| < 1\) and \(x(1) = y(1) = 0\) hold. Since \(\left| \frac{1}{\bar{c}} \right| < 1\) implies \(|c| > 1\), we recognize that this contradicts the assumption \(|c| < 1\) in this case.
    \end{enumerate}

    \item We suppose that the conditions \(\left| \frac{1}{\bar{c}} \right| < 1\) and \(x(1) = y(1) = 0\) are satisfied, which implies \(|c| > 1\). In this case, we find that the eigenvalues are \(\lambda_1 = \frac{1}{\bar{c}}\) and \(\lambda_2 = c\). We introduce a matrix \(T\) defined as follows:

    \begin{align*}
    T = \begin{pmatrix}
    0 & \frac{2a c}{|c|^2 - 1} & 0 & 1 \\
    0 & 1 & 0 & 0 \\
    1 & 0 & -\frac{2a c}{|c|^2 - 1} & 0 \\
    0 & 0 & 1 & 0
    \end{pmatrix}.
    \end{align*}

    Since the determinant of \(T\) is $1$, we conclude that \(T\) is an invertible matrix. The inverse is

    \begin{align*}
    T^{-1} = \begin{pmatrix}
    0 & 0 & 1 & \frac{2a c}{|c|^2 - 1} \\
    0 & 1 & 0 & 0 \\
    0 & 0 & 0 & 1 \\
    1 & -\frac{2a c}{|c|^2 - 1} & 0 & 0
    \end{pmatrix}.
    \end{align*}

    Under this transformation, we verify that \(T^{-1} R T = \text{diag}\left( \frac{1}{\bar{c}}, \frac{1}{\bar{c}}, c, c \right)\). For \(n = 1\), we obtain:

    \begin{align*}
    T^{-1} \psi(1) = z \begin{pmatrix} 0 \\ 0 \\ 0 \\ \bar{c} \end{pmatrix} + w \begin{pmatrix} \frac{2a}{|c|^2 - 1} \\ 0 \\ \bar{c} \\ 0 \end{pmatrix},
    \end{align*}

    and we compute the \(l^2\)-norm of \(T^{-1} \psi\) as:

    \begin{align*}
    \| T^{-1} \psi \|^2 = \sum_{n \in \mathbb{N}} \left( \left| \frac{1}{\bar{c}} \right|^{2n} \big( |x(n)|^2 + |y(n)|^2 \big) + |c|^{2n} \big( |z(n)|^2 + |w(n)|^2 \big) \right).
    \end{align*}

   Given that the product of the absolute values \(|\lambda_1| |\lambda_2|\) is 1 and \(|c| < 1\), \(\| T^{-1} \psi \|^2 < \infty\) holds if and only if either (i) \(z(1) = w(1) = 0\), or (ii) \(x(1) = y(1) = 0\) and \(\left| \frac{1}{\bar{c}} \right| < 1\).

    \begin{enumerate}
        \item We suppose that \(z(1) = w(1) = 0\) holds. Since \(z(1) = z \bar{c} = 0\) and \(w(1) = w \bar{c} = 0\), and given that \(c \neq 0\), we deduce that \(z = 0\) and \(w = 0\). Therefore, we determine that the kernel of \(H^\# - E\) is trivial, i.e., \(\ker(H^\# - E) = \{0\}\).

        \item We suppose that \(|c| < 1\) and \(x(1) = y(1) = 0\) hold. In this case, we recognize that \(|c| < 1\) contradicts the assumption \(|c| > 1\).
    \end{enumerate}
\end{enumerate}
  \end{proof}

\end{document}
